\chapter{Transfer learning y ajuste fino: historia y teoría}\label{cap:historia}

Este capítulo cuenta de dónde viene la idea de transferir lo aprendido y cómo llegó a ser la forma habitual de
entrenar modelos. No es un repaso exhaustivo: elige los trabajos que sirven para entender las decisiones y los
resultados de NOVA. La \cref{fig:linea} resume el recorrido y la \cref{tab:estrategias} ubica cada estrategia de adaptación
en este trabajo.

\section{Una idea más antigua que las redes}

La pregunta de si aprender una cosa ayuda a aprender otra es de la psicología. \textcite{thorndike1901} estudiaron si
mejorar en una función mental mejora la eficiencia en otras funciones. Encontraron que la mejora se traspasa solo en
parte y que depende de lo que las dos tareas tienen en común. La idea de
fondo, que lo aprendido se transfiere mejor cuanto más se parecen las tareas, reaparece casi un siglo después en las
redes neuronales, con otras palabras: la distancia entre dominios.

\section{Las primeras redes que transfieren}

La historia habitual sitúa el comienzo del transfer learning en redes neuronales a principios de los noventa.
\textcite{bozinovski2020} recuerda que el trabajo empezó más de una década antes: hubo investigación en los setenta, con
una primera publicación en 1976, y ya tenía un modelo matemático de la transferencia, una interpretación geométrica y
una medida que distingue entre transferencia positiva, negativa y nula.

En los noventa la idea toma su forma moderna. \textcite{pratt1991} estudiaron cómo usar lo que aprendió una red para
que otra aprenda una tarea relacionada, transfiriendo directamente los pesos, y lo usaron para acelerar el
aprendizaje. \textcite{pratt1993} propuso una forma más fina, la transferencia basada en discriminabilidad: estima con
una medida de información qué tan útil es, en el problema destino, cada hiperplano que definen los pesos de la fuente,
y reescala los pesos transferidos según eso. Las redes que parten así aprenden bastante más rápido que las que parten
al azar. Es, treinta años antes, la intuición detrás del ajuste fino parcial que se discute más abajo: no todas las
partes de una red sirven igual en el dominio nuevo.

\section{Aprender a aprender}

A mediados de los noventa la pregunta cambia de escala. \textcite{thrun1996} preguntó si aprender la tarea número $n$
es más fácil que aprender la primera, y planteó el aprendizaje de por vida: un sistema que acumula conocimiento a lo
largo de muchas tareas y lo usa para aprender las siguientes con menos ejemplos. \textcite{caruana1997} propuso el
aprendizaje multitarea: entrenar a la vez varias tareas relacionadas, con una representación compartida, para que lo
que aprende cada una sirva de sesgo a las demás. El libro editado por \textcite{thrun1998} reunió esta línea bajo el
nombre de \enquote{aprender a aprender}, y \textcite{baxter2000} le dio una base teórica: si las tareas vienen de un
mismo entorno, se puede aprender el sesgo inductivo que conviene a todas, y el número de ejemplos que necesita cada
tarea nueva baja.

En NOVA esta idea aparece en el modelo de pronóstico, que es \emph{compartido entre variables}: un mismo PatchTST
aprende la dinámica diaria del pulso, la presión, el peso y los demás canales, en vez de un modelo por variable.

\section{Qué es transferir, con precisión}\label{sec:definicion}

La definición que se volvió estándar es la de \textcite{pan2010}. Un \emph{dominio} $\mathcal{D} = \{\mathcal{X},
P(X)\}$ es un espacio de entradas con su distribución, y una \emph{tarea} $\mathcal{T} = \{\mathcal{Y}, f(\cdot)\}$ es
un espacio de salidas con la función que se quiere aprender. Dados un dominio y una tarea fuente
$(\mathcal{D}_S, \mathcal{T}_S)$ y un dominio y una tarea destino $(\mathcal{D}_T, \mathcal{T}_T)$, el transfer
learning busca aprender mejor $f_T$ usando lo que se sabe de $\mathcal{D}_S$ y $\mathcal{T}_S$, cuando
$\mathcal{D}_S \neq \mathcal{D}_T$ o $\mathcal{T}_S \neq \mathcal{T}_T$.

Con esa definición, \textcite{pan2010} ordenaron el campo en tres familias según qué cambia y si hay etiquetas en el
destino (inductivo, transductivo y no supervisado) y en cuatro según qué se transfiere: instancias, representaciones,
parámetros o relaciones. Las revisiones posteriores de \textcite{weiss2016}, \textcite{tan2018} y
\textcite{zhuang2021} extendieron esa taxonomía al aprendizaje profundo y a cientos de métodos nuevos.

\begin{clave}
En NOVA, la fuente es la Tierra: la tabla diaria de LifeSnaps y los días de pre-reposo de cada sujeto. El destino es el
reposo en cama, análogo del espacio. La tarea es la misma en ambos (el valor esperado de una variable o sus próximos 7
días), pero cambia la distribución: el peso baja, el pulso sube y la \enquote{normalidad} de cada persona se corre. Es
transfer learning \emph{inductivo con cambio de dominio}, y lo que se transfiere son parámetros (los pesos de PatchTST)
y una línea base personal.
\end{clave}

\section{Cuando cambia la distribución}

El caso más estudiado de cambio de dominio es el \emph{cambio de covariables}: cambia la distribución de las entradas,
$P_S(x) \neq P_T(x)$, pero no la relación entre entrada y salida, $P_S(y \mid x) = P_T(y \mid x)$.
\textcite{shimodaira2000} mostró que en ese caso conviene ponderar la verosimilitud de cada ejemplo de entrenamiento
por el cociente de densidades
\begin{equation}
w(x) = \frac{p_T(x)}{p_S(x)},
\end{equation}
de modo que los ejemplos que se parecen más al destino pesen más.

\textcite{bendavid2010} dieron una cota que dice cuándo un clasificador $h$ entrenado en la fuente puede funcionar en
el destino:
\begin{equation}\label{eq:bendavid}
\varepsilon_T(h) \le \varepsilon_S(h) + \tfrac{1}{2}\, d_{\mathcal{H}\Delta\mathcal{H}}(\mathcal{D}_S, \mathcal{D}_T)
  + \lambda .
\end{equation}
El error en el destino queda acotado por el error en la fuente, más una distancia entre las dos distribuciones que
mide qué tan bien puede distinguirlas la familia de modelos, más $\lambda$, el error del mejor modelo conjunto en
ambos dominios. Si $\lambda$ es grande, ninguna representación compartida sirve y transferir puede empeorar las cosas:
es la \emph{transferencia negativa}. \textcite{ganin2016} convirtieron la cota en un método: entrenar una
representación que prediga bien en la fuente y que, a la vez, un discriminador no pueda usar para distinguir de qué
dominio viene cada ejemplo. Antes, \textcite{dai2007} habían adaptado el \emph{boosting} para repesar instancias de la
fuente según cuánto ayudan en el destino.

En el reposo en cama el cambio no es solo de covariables. El nivel normal del peso de una persona baja con los días,
así que también cambia $P(y \mid x)$: el mismo historial ya no predice el mismo valor. Eso justifica un modelo que
corrija explícitamente el nivel según el día en cama (\cref{sec:rango}).

\section{Pre-entrenar y ajustar fino}

El ajuste fino, tal como se usa hoy, nace con el aprendizaje profundo. \textcite{hinton2006dbn} propusieron entrenar
una red profunda capa por capa, sin etiquetas, y después \enquote{ajustar fino} (\emph{fine-tune}) todos los pesos con
un segundo algoritmo. \textcite{hinton2006science} usaron la misma receta para autoencoders profundos: el
pre-entrenamiento deja los pesos cerca de una buena solución y el ajuste fino con retropropagación termina el trabajo.
\textcite{bengio2007} extendieron el entrenamiento voraz capa por capa a otros tipos de unidades, y
\textcite{erhan2010} explicaron por qué funciona: el pre-entrenamiento actúa como un regularizador que lleva la
optimización a regiones del espacio de parámetros que generalizan mejor.

En paralelo, \textcite{raina2007} propusieron el aprendizaje autodidacta: aprender una representación con datos sin
etiquetar de otra fuente y usarla en una tarea supervisada con pocos ejemplos. \textcite{bengio2012} reunió el
argumento de que las buenas representaciones son justamente las que se transfieren. En notación, el ajuste fino parte
de los parámetros aprendidos en la fuente y sigue optimizando en el destino:
\begin{equation}\label{eq:ajuste}
\theta_T = \arg\min_{\theta}\; \mathcal{L}_T(\theta), \qquad \text{empezando en } \theta = \theta_S ,
\end{equation}
con todos los parámetros libres (ajuste completo) o solo un subconjunto, el resto congelado.

\section{Qué capas se transfieren}

Con las redes convolucionales entrenadas en ImageNet, transferir se volvió la práctica por defecto en visión.
\textcite{donahue2014} mostraron que las activaciones internas de una red ya entrenada, usadas tal cual, sirven como
rasgos genéricos para tareas nuevas, y \textcite{razavian2014} confirmaron que esos rasgos \enquote{sacados del
estante} son una línea de base difícil de superar. \textcite{oquab2014} transfirieron las representaciones intermedias a tareas con pocos
datos, entrenando capas nuevas encima, y \textcite{girshick2014} establecieron la receta que dominó la década: pre-entrenar
supervisado en un conjunto grande y ajustar fino en el dominio específico.

\textcite{yosinski2014} midieron capa por capa qué tan transferibles son los rasgos. Las primeras capas aprenden cosas
generales y las últimas se especializan en la tarea fuente. La transferencia empeora por dos causas: esa
especialización y la coadaptación entre capas vecinas, que se rompe al cortar la red en un punto. Cuanto más distintas
son las tareas, menos sirve transferir, aunque iniciar con rasgos transferidos sigue dando una ventaja que dura después
del ajuste fino.

\section{La normalización también se adapta}\label{sec:bn}

\textcite{ioffe2015} introdujeron la normalización por lotes (BatchNorm) para reducir lo que llamaron cambio interno de
covariables. Cada activación se normaliza con la media y la varianza del lote y se reescala con dos parámetros
entrenables:
\begin{equation}\label{eq:bn}
\hat{x} = \frac{x - \mu_{\mathcal{B}}}{\sqrt{\sigma^2_{\mathcal{B}} + \epsilon}}, \qquad y = \gamma \hat{x} + \beta,
\qquad \mu \leftarrow (1-m)\,\mu + m\,\mu_{\mathcal{B}}, \quad \sigma^2 \leftarrow (1-m)\,\sigma^2 + m\,\sigma^2_{\mathcal{B}}.
\end{equation}
Las medias móviles $\mu$ y $\sigma^2$ no son parámetros que se entrenan con gradientes: son \emph{buffers} que se
actualizan cada vez que el modelo procesa un lote en modo de entrenamiento, y que se usan, fijos, en modo de
evaluación.

Eso convierte a BatchNorm en una vía de adaptación por sí sola. \textcite{li2018adabn} propusieron AdaBN: adaptar un
modelo a otro dominio reemplazando solo esas estadísticas por las del dominio destino, sin tocar los pesos.
\textcite{schneider2020} mostraron que lo mismo mejora la robustez frente a imágenes corrompidas. El punto importa para
este trabajo: el brazo \enquote{solo la cabeza} de NOVA entrena con el modelo en modo de entrenamiento, así que además
de la cabeza adapta las estadísticas de BatchNorm del backbone (\cref{sec:pesos,sec:a1}).

\section{Lo que se olvida al adaptar}

Ajustar una red a un dominio nuevo puede borrar lo que sabía del anterior. \textcite{mccloskey1989} lo llamaron
interferencia catastrófica: al aprender tareas en secuencia, lo nuevo pisa lo viejo. \textcite{kirkpatrick2017}
propusieron frenar ese olvido penalizando el cambio de los pesos importantes para la tarea anterior, con una
importancia estimada por la información de Fisher $F_i$:
\begin{equation}
\mathcal{L}(\theta) = \mathcal{L}_T(\theta) + \sum_i \frac{\lambda}{2}\, F_i\, \bigl(\theta_i - \theta^{*}_{S,i}\bigr)^2 .
\end{equation}
En NOVA el olvido se evita por diseño: el modelo de Tierra no se reemplaza, el adaptado se guarda aparte, y el
corrector del rango esperado se suma a la línea base sin modificarla.

\section{Modelos fundacionales y ajuste eficiente}

El transformer de \textcite{vaswani2017} cambió la escala del pre-entrenamiento. \textcite{howard2018} mostraron con
ULMFiT que un modelo de lenguaje pre-entrenado en texto general se puede ajustar fino a cualquier clasificación de
textos, con técnicas para no destruir lo aprendido: tasas de aprendizaje distintas por capa y descongelar las capas de
a una. \textcite{devlin2019} llevaron la receta a BERT, que se ajusta a cada tarea con una sola capa de salida
adicional. \textcite{bommasani2021} acuñaron el término \emph{modelos fundacionales} para estos modelos entrenados con
datos amplios y adaptables a muchas tareas.

Como ajustar todos los parámetros de un modelo grande es caro, aparecieron formas de ajustar pocos.
\textcite{houlsby2019} insertaron \emph{adaptadores}, módulos pequeños entre las capas congeladas.
\textcite{hu2022} propusieron LoRA: congelar los pesos $W_0$ y aprender una corrección de rango bajo,
\begin{equation}\label{eq:lora}
W = W_0 + BA, \qquad B \in \mathbb{R}^{d \times r},\; A \in \mathbb{R}^{r \times k},\; r \ll \min(d, k).
\end{equation}
En NOVA, LoRA queda para después de la corrida v4, y solo si el ajuste completo sigue ganándole al modelo entrenado
desde cero (\cref{sec:v4}).

\section{¿Cuándo ayuda pre-entrenar?}\label{sec:cuando}

Que el pre-entrenamiento sea la norma no quiere decir que siempre ayude. \textcite{he2019} entrenaron detectores de
objetos desde cero, con más iteraciones, y obtuvieron resultados no peores que con pre-entrenamiento en ImageNet: el
pre-entrenamiento acelera la convergencia, pero no siempre mejora el resultado final. \textcite{raghu2019} encontraron
que en imágenes médicas transferir desde ImageNet aporta poco y que modelos chicos entrenados desde cero rinden
parecido. \textcite{kornblith2019} vieron que en dos conjuntos pequeños de clasificación fina el beneficio es mínimo.

\textcite{neyshabur2020} separaron qué se transfiere: una parte es reutilización de rasgos y otra son estadísticas de
bajo nivel de los datos, que ayudan aunque los datos fuente se desordenen. \textcite{kumar2022} mostraron que el ajuste
completo puede distorsionar los rasgos pre-entrenados y rendir peor que entrenar solo la cabeza cuando el destino se
aleja de la fuente, y \textcite{lee2023} que ajustar solo un subconjunto de capas, elegido según el tipo de cambio de
distribución, puede rendir mejor que ajustarlas todas.

Hay además un problema de medición. \textcite{dodge2020} mostraron que, en el ajuste fino de modelos de lenguaje, la
semilla aleatoria cambia los resultados de forma sustancial, y que la inicialización de la cabeza y el orden de los
datos pesan parecido. \textcite{bouthillier2021} pidieron contar esa variabilidad antes de declarar que un modelo le
gana a otro. Con una sola semilla, una diferencia de milésimas no se puede separar del ruido. Es exactamente la
situación del pronóstico de NOVA (\cref{sec:res-pronostico,sec:h1}).

\section{Transferir en series de tiempo}

En series de tiempo el pre-entrenamiento llegó más tarde. \textcite{fawaz2018} probaron transferir redes
convolucionales entre conjuntos de clasificación de series y vieron que, según la fuente, la transferencia puede
mejorar o empeorar el resultado. \textcite{oreshkin2021} plantearon el pronóstico \enquote{zero-shot}: un modelo
entrenado en muchas series que pronostica series nuevas sin reentrenar. \textcite{zeng2023} pusieron un límite a la
euforia: en pronóstico de largo plazo, un modelo lineal de una capa le gana a varios transformers.

PatchTST \parencite{nie2023}, el modelo base de NOVA, respondió a esa crítica con dos decisiones: cortar la serie en
\emph{parches} de varios pasos, que funcionan como las palabras de un texto, y tratar cada canal por separado con
pesos compartidos (independencia de canales). Después llegaron los modelos fundacionales de series:
Lag-Llama \parencite{rasul2023}, Chronos \parencite{ansari2024}, que escala y cuantiza la serie en fichas y la modela
con la familia T5, TimesFM \parencite{das2024}, Moirai \parencite{woo2024} y MOMENT \parencite{goswami2024}.
Chronos-2 \parencite{ansari2025}, la referencia de este trabajo, pronostica sin ajuste series univariadas,
multivariadas y con covariables, con atención entre las series de un mismo grupo.

Para comparar pronósticos entre variables de unidades distintas hacen falta errores relativos.
\textcite{hyndman2006} discutieron las medidas escaladas y relativas a un pronóstico ingenuo. Este trabajo usa el MAE
dividido por el MAE de la persistencia (repetir el último valor), en esa línea.

\section{Línea de tiempo}

La \cref{fig:linea} junta en una sola línea los trabajos de este capítulo, desde la psicología de
principios del siglo~\textsc{xx} hasta los modelos fundacionales de series de tiempo, y marca dónde entra NOVA. Los colores
separan las ideas y la teoría, el pre-entrenamiento y el ajuste fino, y las series de tiempo.

\begin{figure}[p]
\centering
\newcommand*{\hito}[4]{%
  \node[anchor=east, font=\small\titularfont\bfseries\color{azulnoche}] at (-0.32,#1) {#2};
  \node[circle, inner sep=0pt, minimum size=6.5pt, fill=#3, draw=white, line width=1.1pt] at (0,#1) {};
  \node[anchor=west, font=\footnotesize, text width=12.3cm, align=left, inner sep=1pt] at (0.32,#1) {#4};}
\begin{tikzpicture}[y=-7.65mm]
  \draw[reja, line width=1.8pt] (0,-0.5) -- (0,26.5);
  \hito{0}{1901}{tinta3}{\citeauthor{thorndike1901}: mejorar en una tarea mental se traspasa a otras solo en parte.}
  \hito{1}{1976}{tinta3}{Primera publicación sobre transferencia en redes neuronales, con una medida de transferencia positiva, negativa y nula \parencite{bozinovski2020}.}
  \hito{2}{1989}{tinta3}{\citeauthor{mccloskey1989}: interferencia catastrófica al aprender en secuencia.}
  \hito{3}{1991}{tinta3}{\citeauthor{pratt1991}: transferir los pesos de una red como punto de partida de otra.}
  \hito{4}{1992}{tinta3}{\citeauthor{pratt1993}: transferencia basada en discriminabilidad, que reescala los pesos transferidos.}
  \hito{5}{1995}{tinta3}{\citeauthor{thrun1996}: aprendizaje de por vida; ¿es más fácil aprender lo enésimo?}
  \hito{6}{1997}{tinta3}{\citeauthor{caruana1997}: aprendizaje multitarea con representación compartida.}
  \hito{7}{2000}{tinta3}{\citeauthor{baxter2000}: aprender el sesgo inductivo. \citeauthor{shimodaira2000}: cambio de covariables.}
  \hito{8}{2006}{serieA}{\citeauthor{hinton2006dbn}: pre-entrenar capa por capa y después ajustar fino.}
  \hito{9}{2007}{serieA}{\citeauthor{raina2007}: aprendizaje autodidacta. \citeauthor{dai2007}: TrAdaBoost.}
  \hito{10}{2009}{tinta3}{\citeauthor{bendavid2010}: cota del error en el dominio destino.}
  \hito{11}{2010}{tinta3}{\citeauthor{pan2010}: definición y taxonomía del transfer learning. \citeauthor{erhan2010}: por qué ayuda pre-entrenar.}
  \hito{12}{2014}{serieA}{\citeauthor{yosinski2014}: rasgos generales abajo, específicos arriba. \citeauthor{girshick2014}: pre-entrenar y ajustar fino en visión.}
  \hito{13}{2015}{serieA}{\citeauthor{ioffe2015}: BatchNorm.}
  \hito{14}{2016}{tinta3}{\citeauthor{ganin2016}: adaptación de dominio adversaria.}
  \hito{15}{2017}{serieA}{\citeauthor{vaswani2017}: el transformer. \citeauthor{kirkpatrick2017}: frenar el olvido con EWC.}
  \hito{16}{2018}{serieA}{\citeauthor{howard2018}: ULMFiT. \citeauthor{li2018adabn}: adaptar solo las estadísticas de BatchNorm. \citeauthor{fawaz2018}: transferir en series de tiempo.}
  \hito{17}{2019}{serieA}{\citeauthor{devlin2019}: BERT. \citeauthor{houlsby2019}: adaptadores. \citeauthor{he2019}: entrenar desde cero no siempre pierde.}
  \hito{18}{2020}{serieA}{\citeauthor{neyshabur2020}: qué se transfiere. \citeauthor{dodge2020}: el peso de la semilla en el ajuste fino.}
  \hito{19}{2021}{serieA}{\citeauthor{hu2022}: LoRA. \citeauthor{bommasani2021}: modelos fundacionales.}
  \hito{20}{2022}{serieC}{\citeauthor{nie2023}: PatchTST. \citeauthor{kumar2022}: el ajuste fino puede distorsionar rasgos. \citeauthor{lee2023}: ajuste quirúrgico.}
  \hito{21}{2023}{serieC}{\citeauthor{zeng2023}: un modelo lineal le gana a los transformers. \citeauthor{rasul2023}: Lag-Llama. \citeauthor{das2024}: TimesFM.}
  \hito{22}{2024}{serieC}{\citeauthor{ansari2024}: Chronos. \citeauthor{woo2024}: Moirai. \citeauthor{goswami2024}: MOMENT.}
  \hito{23}{2025}{serieC}{\citeauthor{ansari2025}: Chronos-2, univariado, multivariado y con covariables.}
  \hito{24.4}{2026}{serieB}{\textbf{NOVA:} línea base de Tierra + corrector de otros sujetos + días propios, y PatchTST Tierra → reposo en cama.}
  % leyenda
  \begin{scope}[shift={(0.32,26.3)}]
    \node[circle, inner sep=0pt, minimum size=6pt, fill=tinta3] at (0,0) {};
    \node[anchor=west, font=\scriptsize\sffamily\color{tinta2}] at (0.15,0) {ideas y teoría};
    \node[circle, inner sep=0pt, minimum size=6pt, fill=serieA] at (2.9,0) {};
    \node[anchor=west, font=\scriptsize\sffamily\color{tinta2}] at (3.05,0) {pre-entrenamiento y ajuste fino};
    \node[circle, inner sep=0pt, minimum size=6pt, fill=serieC] at (7.6,0) {};
    \node[anchor=west, font=\scriptsize\sffamily\color{tinta2}] at (7.75,0) {series de tiempo};
  \end{scope}
\end{tikzpicture}
\caption[Línea de tiempo del transfer learning y del ajuste fino]{Línea de tiempo del transfer learning y del ajuste
fino, con los trabajos que usa este informe. Los años son los de la versión citada en la bibliografía.}
\label{fig:linea}
\end{figure}

\section{Dónde cae este trabajo}

La \cref{tab:estrategias} ordena las estrategias de adaptación de más simple a más compleja, como pide el plan del
proyecto, y dice cuántos parámetros cambia cada una en NOVA y si está construida.

\begin{table}[htbp]
\caption[Estrategias de adaptación y su lugar en este trabajo]{Estrategias de adaptación, de la más simple a la más
compleja, y su lugar en este trabajo. El PatchTST tiene \cifra{modelo.parametros} parámetros entrenables.}
\label{tab:estrategias}
\small
\begin{tabularx}{\textwidth}{@{}>{\raggedright\arraybackslash}p{3.2cm} >{\raggedright\arraybackslash}X
  >{\raggedright\arraybackslash}p{2.7cm} >{\raggedright\arraybackslash}p{2.5cm}@{}}
\toprule
Estrategia & Qué se ajusta en NOVA & Cuánto cambia & Estado \\
\midrule
Corrector residual & Tres coeficientes ridge por variable (rango esperado), más la media contraída de los días
propios & 3 por variable & Construido; es el resultado fuerte \\
Sin adaptar & Nada: el modelo de Tierra o Chronos-2 se usan tal cual & 0 & Construido, como referencia \\
Solo la cabeza & La proyección final del PatchTST & \cifra{modelo.cabeza} parámetros (\cifra{modelo.cabeza.pct}\pct)
& Construido, pero también mueve BatchNorm \\
Estadísticas de BatchNorm & Medias y varianzas móviles del backbone, sin gradientes (AdaBN) & \cifra{pesos.bn}
tensores & Ocurre dentro de \enquote{solo la cabeza} \\
Ajuste completo & Todos los parámetros, con una tasa de aprendizaje baja & \cifra{modelo.parametros} parámetros &
Construido \\
Adaptadores o LoRA & Matrices de rango bajo sobre los pesos congelados & Según el rango & Diseñado; solo después de
la v4 \\
Desde cero & Todos los parámetros, sin pre-entrenamiento & \cifra{modelo.parametros} parámetros & Construido, como
control \\
\bottomrule
\end{tabularx}
\end{table}
